\documentclass[11pt]{article}
\usepackage[margin=1in]{geometry}
\usepackage{amsmath,amssymb,amsthm}
\usepackage{booktabs}
\usepackage{microtype}
\usepackage{hyperref}
\usepackage{xcolor}
\usepackage{listings}

\lstset{
  language=C++,
  basicstyle=\ttfamily\small,
  keywordstyle=\color{blue},
  commentstyle=\color{gray},
  breaklines=true,
  frame=single,
}

\title{\texttt{EOM::ArnoldiSolve} --- Algorithm Explained}
\author{IMSRG project notes}
\date{\today}

\begin{document}
\maketitle
\tableofcontents
\bigskip

This note explains the algorithm implemented in
\texttt{src/EOM.cc} (\texttt{EOM::ArnoldiSolve}), which performs a
Rayleigh--Ritz eigenvalue calculation for a Hamiltonian split as
\begin{equation}
  H \;=\; H_1 + H_2,
\end{equation}
where
\begin{itemize}
  \item $H_1 v$ can be applied directly (cheap one-sided action).
  \item $H_2 v$ cannot be formed; only \emph{diagonal} expectation
        values $\langle u \mid H_2 \mid u \rangle$ are affordable.
  \item $H = H_1 + H_2$ is symmetric in the metric defined by
        \texttt{ComputeNorm}$(a,b) = \langle a \mid b \rangle$,
        \textbf{but $H_1$ and $H_2$ individually are not symmetric}
        in that metric.
\end{itemize}
The key observation is that even though $H_2$ is never applied to a
vector, its off-diagonal matrix elements
$\langle v_i \mid H_2 \mid v_j \rangle$ between two \emph{known}
vectors can still be recovered from three diagonal expectation values
via a polarization identity.

% ----------------------------------------------------------------
\section{Subspace built from $H_1$ only}
% ----------------------------------------------------------------

We build a metric-orthonormal Krylov-like basis
$\{v_0, v_1, \dots, v_{m-1}\}$ by applying only $H_1$:

\begin{enumerate}
  \item Normalize the initial vector:
        $v_0 = v_i \big/ \sqrt{\langle v_i \mid v_i \rangle}$.
  \item At step $j$, compute and cache
        $\tilde{h}_1 v_j \equiv H_1 v_j$
        (\texttt{HtcMultiref(Hs, $v_j$)}).
  \item Project onto the physical subspace (\texttt{ProjectOprator}).
  \item Double-pass classical Gram--Schmidt against
        $\{v_0, \dots, v_j\}$ in the metric \texttt{ComputeNorm},
        projecting between passes.
  \item Set
        \begin{equation}
          v_{j+1} \;=\; \frac{w}{\sqrt{\langle w \mid w \rangle}}.
        \end{equation}
\end{enumerate}

The resulting subspace is
\begin{equation}
  \mathcal{K}_m \;=\;
  \operatorname{span}\{v_0,\; H_1 v_0,\; H_1^2 v_0,\; \dots\},
\end{equation}
the Krylov space of $H_1$ \emph{alone} --- not of the full $H$.
This is precisely why we never need to apply $H_2$.

% ----------------------------------------------------------------
\section{The polarization trick for the $H_2$ off-diagonals}
% ----------------------------------------------------------------

For any operator $A$ that is \emph{symmetric} in the inner product,
\begin{equation}
  \langle u \mid A \mid v \rangle
  \;=\;
  \tfrac{1}{2}\Bigl(
    \langle u+v \mid A \mid u+v \rangle
    - \langle u \mid A \mid u \rangle
    - \langle v \mid A \mid v \rangle
  \Bigr).
\end{equation}

In our case $H_2$ is \textbf{not} symmetric, but its symmetrized
part is what appears in $H = H_1 + H_2$.  Define
\begin{equation}
  \langle u \mid H_2 \mid v \rangle_{\mathrm{sym}}
  \;\equiv\;
  \tfrac{1}{2}\bigl(\langle u \mid H_2 \mid v \rangle
                   + \langle v \mid H_2 \mid u \rangle\bigr),
\end{equation}
and similarly for $H_1$.  Polarization recovers exactly the
symmetrized piece:
\begin{equation}
  \langle v_i \mid H_2 \mid v_j \rangle_{\mathrm{sym}}
  \;=\;
  \tfrac{1}{2}\Bigl(
    \underbrace{\langle v_i+v_j \mid H_2 \mid v_i+v_j \rangle}_{\text{cross term}}
    - \underbrace{\langle v_i \mid H_2 \mid v_i \rangle}_{h_2^{\mathrm{diag}}[i]}
    - \underbrace{\langle v_j \mid H_2 \mid v_j \rangle}_{h_2^{\mathrm{diag}}[j]}
  \Bigr).
\end{equation}
Every term on the right is a \emph{diagonal} expectation value, which
is affordable by assumption (\texttt{DcomMultiref}).

We cache $h_2^{\mathrm{diag}}[k] = \langle v_k \mid H_2 \mid v_k \rangle$
for each basis vector the moment it is added.  Consequently, each new
step $j$ requires only \textbf{one new cross term}
$\langle v_i+v_j \mid H_2 \mid v_i+v_j \rangle$ per off-diagonal
entry $(i,j)$ with $i < j$.

% ----------------------------------------------------------------
\section{The $H_1$ off-diagonals}
% ----------------------------------------------------------------

For $H_1$ we do have the one-sided action.  Its \emph{asymmetric}
matrix elements are
\begin{equation}
  \langle v_i \mid H_1 \mid v_j \rangle
  = \langle v_i \mid H_1 v_j \rangle,
  \qquad
  \langle v_j \mid H_1 \mid v_i \rangle
  = \langle v_j \mid H_1 v_i \rangle,
\end{equation}
both computed via \texttt{ComputeNorm} using the cached $H_1 v_k$
from every previous step.  We take the symmetric part:
\begin{equation}
  \langle v_i \mid H_1 \mid v_j \rangle_{\mathrm{sym}}
  \;=\;
  \tfrac{1}{2}\bigl(
    \langle v_i \mid H_1 v_j \rangle
    + \langle v_j \mid H_1 v_i \rangle
  \bigr).
\end{equation}

% ----------------------------------------------------------------
\section{Assembling the projected matrix}
% ----------------------------------------------------------------

With both pieces symmetrized, the projected Hamiltonian matrix element
is
\begin{equation}
  \boxed{
  H_{ij}
  \;=\;
  \langle v_i \mid H_1 \mid v_j \rangle_{\mathrm{sym}}
  \;+\;
  \langle v_i \mid H_2 \mid v_j \rangle_{\mathrm{sym}}
  }
\end{equation}
stored in \texttt{hall(i,j) = hall(j,i)}.

\subsection*{Diagonal short-circuit ($i = j$)}

For $i = j$ the polarization formula collapses analytically:
\begin{equation}
  \tfrac{1}{2}\bigl(
    \langle 2v \mid H_2 \mid 2v \rangle
    - 2\langle v \mid H_2 \mid v \rangle
  \bigr)
  \;=\;
  \tfrac{1}{2}(4 - 2)\langle v \mid H_2 \mid v \rangle
  \;=\;
  \langle v \mid H_2 \mid v \rangle
  \;=\;
  h_2^{\mathrm{diag}}[j].
\end{equation}
On the diagonal we therefore
\begin{itemize}
  \item \textbf{skip} the extra \texttt{DcomMultiref(Hs, $2v_j$)} call
        (saves one expensive evaluation per step), and
  \item \textbf{avoid} the catastrophic cancellation
        $4x - 2x \to x$ that would otherwise lose one decimal digit of
        precision when $\langle v \mid H_2 \mid v \rangle$ is large.
\end{itemize}
Likewise for $H_1$, the diagonal is just
$\langle v_j \mid H_1 v_j \rangle$ directly.

% ----------------------------------------------------------------
\section{Why explicit symmetrization matters}
% ----------------------------------------------------------------

In general $\langle v_i \mid H_1 \mid v_j \rangle \neq
\langle v_j \mid H_1 \mid v_i \rangle$, and the same holds for $H_2$.
The combined matrix element is symmetric \emph{because the antisymmetric
parts of $H_1$ and $H_2$ exactly cancel}.  The polarization trick
delivers the \textbf{symmetrized} $H_2$ piece directly, and we
explicitly symmetrize the $H_1$ piece before adding.  This guarantees
that \texttt{hall(i,j)} is symmetric \textbf{even if} the antisymmetric
parts do not cancel perfectly in finite-precision arithmetic.

\subsection*{Asymmetry diagnostic}

The code tracks, for every off-diagonal pair $(i, j)$,
\begin{equation}
  r_{ij}
  \;=\;
  \frac{
    \bigl|\langle v_i \mid H_1 v_j \rangle
         - \langle v_j \mid H_1 v_i \rangle\bigr|
  }{
    2\,|H_{ij}|
  }
\end{equation}
and reports $\max_{i < j} r_{ij}$ at the end of the run.

A large value of $r_{ij}$ is \emph{allowed} when $H_2$ carries a
matching antisymmetric part; it only signals a problem when the
$H_1/H_2$ split in the calling code is internally inconsistent.
A warning is printed when $\max r_{ij} > 10^{-3}$.

% ----------------------------------------------------------------
\section{Diagonalization with explicit symmetrization}
% ----------------------------------------------------------------

Even after writing \texttt{hall(i,j) = hall(j,i) = value}, repeated
reads/re-writes leave residual bit-level asymmetry that can confuse
\texttt{arma::eig\_sym}.  Every diagonalization is therefore routed
through the helper lambda

\begin{lstlisting}
auto eigensolve_sub = [&hall](int dim, arma::vec &ev, arma::mat &evec)
{
    arma::mat sub = hall.submat(0, 0, dim - 1, dim - 1);
    sub = arma::symmatu(0.5 * (sub + sub.t()));
    arma::eig_sym(ev, evec, sub);
};
\end{lstlisting}

so \texttt{eig\_sym} always receives an exactly symmetric matrix.

% ----------------------------------------------------------------
\section{Convergence and breakdown handling}
% ----------------------------------------------------------------

After each new basis vector, if $j + 1 \ge \texttt{min\_iter}$:
\begin{enumerate}
  \item Diagonalize the leading $(j{+}1)\times(j{+}1)$ block of
        \texttt{hall} via \texttt{eigensolve\_sub}.
  \item Take the lowest \texttt{state\_want} eigenvalues as the current
        Ritz estimates $\mathbf{e}$.
  \item Compare to the previous step's $\mathbf{e}_{\mathrm{prev}}$.
        Declare convergence and stop if
        \begin{equation}
          \max_k |e_k - e_k^{\mathrm{prev}}| < \texttt{tol}
          \quad\text{or}\quad
          \frac{\max_k |e_k - e_k^{\mathrm{prev}}|}{\max(1,\,\max_k|e_k|)}
          < \texttt{tol}.
        \end{equation}
\end{enumerate}

Three breakdown modes can stop the iteration early; each one solves
the current subspace via \texttt{eigensolve\_sub} and returns:

\begin{center}
\begin{tabular}{@{}ll@{}}
  \toprule
  Condition & Meaning \\
  \midrule
  $|b_j| < \texttt{null\_tol}\cdot\texttt{cn0}$
    & New vector exhausted by projection (null space). \\
  $|b_j| < \texttt{bj\_tol}$
    & Exact breakdown: $w \approx 0$. \\
  $b_j < 0$
    & Indefinite metric: \texttt{ComputeNorm} returned a negative norm. \\
  \bottomrule
\end{tabular}
\end{center}

A final eigensolve over the entire accumulated subspace is performed
unconditionally before returning, so the result always reflects the
full basis that was built.

% ----------------------------------------------------------------
\section{Variational consistency check}
% ----------------------------------------------------------------

The function \texttt{EOM::ExpectationValue(Psi)} computes
\begin{equation}
  \mathcal{E}(\Psi)
  \;=\;
  \frac{\langle\Psi \mid H_1 \mid \Psi\rangle
       + \langle\Psi \mid H_2 \mid \Psi\rangle}
       {\langle\Psi \mid \Psi\rangle}
\end{equation}
\emph{independently} of the Krylov bookkeeping --- it uses only
\texttt{HtcMultiref}, \texttt{DcomMultiref}, and
\texttt{ComputeNorm} on the passed vector.

At every 5th Arnoldi step the code builds each current Ritz vector
\begin{equation}
  |\Psi_k\rangle \;=\; \sum_{m=0}^{n_b-1} c_{mk}\, v_m
\end{equation}
and prints
\begin{itemize}
  \item the Ritz value $E_k$ from diagonalizing \texttt{hall},
  \item the direct expectation $\mathcal{E}(\Psi_k)$, and
  \item the difference $\mathcal{E}(\Psi_k) - E_k$.
\end{itemize}
The two numbers \emph{must} agree up to machine precision provided the
basis is metric-orthonormal and \texttt{hall} is assembled correctly.
A growing discrepancy reveals either loss of orthogonality or leakage
of unphysical components into the basis.

% ----------------------------------------------------------------
\section{Ritz vectors}
% ----------------------------------------------------------------

For each eigenvalue $E_k$ with small-space eigenvector
$\mathbf{c}_k = (c_{0k}, c_{1k}, \dots)^T$ the corresponding Ritz
vector in the full operator space is
\begin{equation}
  |\Psi_k\rangle \;=\; \sum_{m=0}^{n_b-1} c_{mk}\, v_m,
\end{equation}
assembled in the loop that fills \texttt{ritz\_vecs}.

% ----------------------------------------------------------------
\section{What the algorithm is not}
% ----------------------------------------------------------------

\begin{itemize}
  \item It is \textbf{not} a standard Lanczos or Arnoldi iteration
        for $H = H_1 + H_2$, because the basis is generated solely
        from $H_1$.  It is a \emph{Galerkin / Rayleigh--Ritz projection}
        of the full $H$ onto $\mathcal{K}_m(H_1, v_0)$.
  \item Convergence of the projected eigenvalues (the
        $\delta < \texttt{tol}$ test) means the small matrix has
        stabilized, \textbf{not} that the Ritz vectors are true
        eigenvectors of $H$.  No $H_2$-applied residual is available,
        so a full residual check on $H$ is not possible.
  \item The method is well-suited to the regime where $H_2$ is small
        or strongly aligned with $H_1$-Krylov directions.  Large,
        structurally different $H_2$ contributions can be missed by
        the subspace.
\end{itemize}

% ----------------------------------------------------------------
\section{Cost summary per Arnoldi step $j$}
% ----------------------------------------------------------------

\begin{center}
\begin{tabular}{@{}lll@{}}
  \toprule
  Quantity & Count at step $j$ & Cost class \\
  \midrule
  $H_1 v_j$ &
    1 &
    one-sided $H_1$ action \\
  $\langle v_i \mid H_1 v_j\rangle,\ \langle v_j \mid H_1 v_i\rangle$ &
    $2j$ inner products &
    cheap \\
  $\langle v_i{+}v_j \mid H_2 \mid v_i{+}v_j\rangle$, $i < j$ &
    $j$ &
    one diagonal $H_2$ expectation each \\
  $\langle v_{j+1} \mid H_2 \mid v_{j+1}\rangle$ (cache) &
    1 &
    one diagonal $H_2$ expectation \\
  Diagonal $\langle v_j \mid H_2 \mid v_j\rangle$ &
    0 \ (uses cache) &
    --- \\
  GS reorthogonalization &
    $2(j{+}1)$ inner products + projections &
    cheap \\
  Eigensolve of $(j{+}1)\times(j{+}1)$ &
    1 &
    tiny dense \\
  \bottomrule
\end{tabular}
\end{center}

The total number of expensive $H_2$ diagonal evaluations after $m$
steps is
\begin{equation}
  \binom{m}{2} + m \;=\; \tfrac{m(m+1)}{2} \;=\; \mathcal{O}(m^2),
\end{equation}
the same scaling as the off-diagonal block we fill --- and with no
$H_2 v$ ever required.

\end{document}
